% TOP_PROBLEM: {"id": 30006929, "title": "Chen-Yang Volume Conjecture for Turaev-Viro Invariants", "category_id": 7, "category": {"id": 7, "name": "topology", "display_name": "Topology", "description": "Properties preserved under continuous deformations.", "slug": "topology", "order_index": 7, "created_at": "2026-07-31T15:26:25.671Z"}, "status": "open"}
% FIRST_POSTED: 2026-09-25
% ENTRY_KIND: statement_only
% !TeX program = lualatex
% Definition and sources only; this file is not an LLM solution attempt.
\documentclass[11pt]{article}
\usepackage[margin=25mm]{geometry}
\usepackage{fontspec}
\setmainfont{Latin Modern Roman}
\usepackage{amsmath,amssymb,mathtools}
\usepackage{unicode-math}
\setmathfont{Latin Modern Math}
\usepackage{xurl,hyperref}
\hypersetup{colorlinks=true,urlcolor=blue}
\setcounter{secnumdepth}{0}
\setlength{\parindent}{0pt}
\setlength{\parskip}{5pt}
\setlength{\emergencystretch}{3em}
\begin{document}
\section*{\texorpdfstring{Chen-Yang Volume Conjecture for Turaev-Viro Invariants}{Chen-Yang Volume Conjecture for Turaev-Viro Invariants}}\label{30006929}
\subsection{Definitions and mathematical statement}
Let $M$ be a finite-volume hyperbolic three-manifold in the closed, cusped, or totally geodesic boundary class specified in Chen--Yang's Conjecture 1.1. Its hyperbolic volume $\operatorname{Vol}(M)$ uses sectional curvature $-1$. For odd integers $r$, $\operatorname{TV}_r(M;q)$ is the Turaev--Viro quantum invariant: a finite sum over admissible labels of a triangulation using quantum dimensions and quantum $6j$-symbols. For cusps and boundary, use precisely the compact-core extension and normalization in the source.

The assertion is
\[
 \lim_{\substack{r\to\infty\\r\text{ odd}}}
 \frac{2\pi}{r}\log\operatorname{TV}_r(M;e^{2\pi i/r})
 =\operatorname{Vol}(M).
\]
The choice of root of unity is part of the problem. The logarithm and invariant normalization are those of the stated source; changing roots or replacing this invariant by a different quantum invariant changes the target. This is an asymptotic growth statement, not equality at finite level.
\par\smallskip\textit{Formulation and concept references:} \href{https://arxiv.org/html/1503.02547}{[S1]}.

\subsection{Short English statement}
Does the exponential growth rate of these Turaev--Viro invariants recover hyperbolic volume for every manifold in the stated class?
\subsection{Sources}
\begin{itemize}
\item[S1]Qingtao Chen and Tian Yang, Volume conjectures for the Reshetikhin-Turaev and the Turaev-Viro invariants, v4 (2018).
\url{https://arxiv.org/html/1503.02547}.
\textit{Formulation source.}
\end{itemize}
\end{document}
